% Part: second-order-logic
% Chapter: sol-and-sets
% Section: cardinalities

\documentclass[../../../include/open-logic-section]{subfiles}

\begin{document}

\olfileid[fr]{sol}{set}{crd}

\olsection{Cardinalités des ensembles}

\begin{explain}
De même que nous pouvons exprimer qu'un
!!{domain} est fini ou infini, !!{enumerable}
ou !!{nonenumerable}, nous pouvons définir
ces propriétés pour un sous-ensemble de
$\Domain{M}$.
\end{explain}

\begin{prop}
La !!{formula} $\fn{Inf}(X)\ident$
\begin{multline*}
\lexists[u][(\lforall[x][(X(x)\lif X(u(x)))]\land{}\\
\lforall[x][\lforall[y][((X(x)\land X(y)\land
\eq[u(x)][u(y)])\lif\eq[x][y])]]\land{}\\
\lexists[y][(X(y)\land\lforall[x][(X(x)
\lif\eq/[y][u(x)])])])]
\end{multline*}
est satisfaite sous une assignation~$s$ si et
seulement si $s(X)$ est infini au sens de
Dedekind. En présence du choix dénombrable,
cela équivaut à l'infinitude ordinaire.
\end{prop}

\begin{prop}
La !!{formula} $\fn{Count}(X)\ident$
\begin{multline*}
\lnot\lexists[x][X(x)]\lor{}\\
\lexists[z][\lexists[u][(X(z)\land
\lforall[x][(X(x)\lif X(u(x)))]\land{}\\
\lforall[Y][((Y\subseteq X\land Y(z)\land
\lforall[x][(Y(x)\lif Y(u(x)))])\lif X=Y)]])]
\end{multline*}
est satisfaite sous une assignation~$s$ si et
seulement si $s(X)$ est !!{enumerable}.
Le premier terme de la disjonction couvre
également l'ensemble vide.
\end{prop}

Le théorème de Cantor montre qu'il existe des
ensembles !!{nonenumerable}. Il existe même
une infinité de niveaux distincts parmi les
cardinalités infinies. La théorie des ensembles
développe une arithmétique de ces tailles et
associe des nombres cardinaux infinis aux
ensembles. Les entiers naturels servent de
cardinaux pour les ensembles finis.
La cardinalité des ensembles !!{denumerable}s
est le premier cardinal infini, noté~$\aleph_0$
(« aleph zéro »). Le cardinal infini suivant,
$\aleph_1$, est le plus petit cardinal non
dénombrable, dans le cadre usuel avec choix.

Nous pouvons exprimer « $X$ est de cardinal
$\aleph_0$ » par
\[
\fn{Aleph}_0(X)\liff
\fn{Inf}(X)\land\fn{Count}(X).
\]
Sous la sémantique standard et avec l'axiome
du choix, $X$ est de cardinal~$\aleph_1$ si
et seulement s'il n'est pas !!{enumerable}
et si chacun de ses sous-ensembles est soit
!!{enumerable}, soit équipotent à~$X$.
Nous pouvons donc poser
\[
\fn{Aleph}_1(X)\ident
\lnot\fn{Count}(X)\land
\lforall[Y][(Y\subseteq X\lif
(\fn{Count}(Y)\lor\cardeq{Y}{X}))].
\]
On peut procéder de manière analogue pour
$\aleph_2$ et les cardinaux suivants.

Une taille présente un intérêt particulier :
la cardinalité du \emph{continu}. C'est celle
de $\Pow{\Nat}$ ou, de façon équivalente,
de~$\Real$. La logique du second ordre peut
aussi exprimer qu'un ensemble a cette taille,
mais la construction demande davantage de
travail.

Dans $\fn{Inf}(X)$, les trois clauses disent que la
restriction de $u$ à~$X$ envoie $X$ dans lui-même, est injective
et omet au moins un élément de~$X$. Réciproquement, toute
injection non surjective de $X$ dans lui-même se prolonge
à une fonction sur le domaine en attribuant une valeur fixe
à tous les arguments hors de~$X$.

Pour $X$ non vide, $\fn{Count}(X)$ demande qu'il soit engendré
par itération de $u$ à partir de~$z$. En effet, l'orbite
$\{z,u(z),u(u(z)),\ldots\}$ est un sous-ensemble de $X$
stable par $u$ et contenant~$z$. La quantification sur~$Y$
impose que cette orbite soit~$X$. Réciproquement, toute
énumération sans répétition de $X$ définit un tel successeur
sur la liste. Si elle est finie, le successeur du dernier
élément est ce dernier élément.

La formule $\fn{Aleph}_1(X)$ s'appuie sur la comparaison
des cardinaux sous l'axiome du choix. Un ensemble de cardinal
$\aleph_1$ ne possède que des sous-ensembles au plus
dénombrables ou de cardinal $\aleph_1$. Si le cardinal
de~$X$ était strictement plus grand que $\aleph_1$, un
sous-ensemble de cardinal $\aleph_1$ contredirait la condition
sur~$Y$.

\end{document}
